\tcbset{
    every box/.style={
        colframe=black,colback=white,
        size=minimal,
        left=2mm,right=2mm,top=2.5mm,bottom=2mm,
        toptitle=1mm,bottomtitle=1mm,
        boxrule=0.8pt,arc=0pt,
    },
    defstyle/.style={fonttitle=\bfseries\upshape, fontupper=\slshape,
        arc=0mm, colback=blue!5!white,colframe=blue!75!black},
    theostyle/.style={fonttitle=\bfseries\upshape, fontupper=\slshape,
        colback=red!10!white,colframe=red!75!black},
    docstyle/.style={theorem style=standard,enhanced,
        colback=docbgcolor,colframe=docfgcolor,coltitle=white,
        fonttitle=\bfseries\upshape,arc=0pt,
        separator sign={\ $\blacktriangleright$},
        drop fuzzy shadow,boxrule=0.4pt},
    donneesstyle/.style={theorem style=standard,enhanced,
        colback=donneesbgcolor,colframe=donneesfgcolor,coltitle=white,
        fonttitle=\bfseries\upshape,arc=0pt,
        separator sign={\ $\blacktriangleright$},
        drop fuzzy shadow,boxrule=0.4pt},
    partstyle/.style={enhanced,bicolor,sidebyside,lefthand width=1cm,
        sharp corners,boxrule=.4pt,
        colframe=partfgcolor,colback=partbgcolor,colbacklower=white,
        fontupper=\LARGE\bfseries,fontlower=\Large\bfseries,halign=center,
        drop fuzzy shadow},
    reponsestyle/.style={enhanced,size=small,sharp corners,
        frame hidden,
        colback=repbgcolor,
        boxrule=0pt,fonttitle=\bfseries,titlerule=0pt,
        colbacktitle=repfgcolor,
        %titlerule hidden,
        title style={left color=repfgcolor,right color=repbgcolor},
        drop fuzzy shadow},
    attentionstyle/.style={enhanced,sharp corners,
        colback=attentionbgcolor,colframe=attentionfgcolor,
        top=1.5mm,bottom=1.5mm,left=0mm,right=0mm,
        boxrule=0pt,leftrule=1cm,
        drop fuzzy shadow,
        overlay={
            \begin{tcbclipframe}
                \node[font=\bfseries,isosceles triangle,inner sep=0pt,
                    isosceles triangle stretches,shape border rotate=90,
                    minimum width=6mm,minimum height=5.5mm,
                    rounded corners=2pt,very thick,
                    text=attentionbgcolor,draw=attentionbgcolor
                ] at ([xshift=5mm]frame.west) {!};
            \end{tcbclipframe}
        }},
    appstyle/.style={theorem style=standard,enhanced,
        colback=appbgcolor,colframe=appfgcolor,coltitle=white,
        fonttitle=\bfseries\upshape,arc=0pt,
        separator sign={\ $\blacktriangleright$},
        drop fuzzy shadow,boxrule=0.4pt,breakable,size=small},
    tcb1style/.style={
        enhanced,
        colframe=secondarycolor,colback=tertiarycolor!95!black,
        top=5mm,boxrule=0.8pt,sharp corners,rounded corners=east,
        varwidth boxed title=0.7\linewidth,
        attach boxed title to top left={%
            yshift=-\tcboxedtitleheight/2-0.25mm,xshift=-5mm},
        boxed title style={%
            sharp corners,rounded corners=east,frame hidden,boxrule=0pt,
        },
        %underlay boxed title={
        %    \draw[red] (title.north west) rectangle (title.south east);
        %},
        colbacktitle=secondarycolor,fonttitle=\bfseries,title={Title},
        %before upper={%
        %},
    },
    LeftRightBorderBoxStyle/.style={%
        enhanced,skin=enhancedmiddle,
        frame hidden,top=0mm,bottom=0.5mm,boxsep=1mm,
        colback=PageBackgroundColor,sharp corners,
        borderline={2mm}{0mm}{#1},
        fuzzy shadow={0mm}{0mm}{-0.7mm}{0.2mm}{black!50},
    },
}


\newtcbox{\qitem}[3]{enhanced,on line,size=minimal,
    auto outer arc,octogon arc,
    colback=#1,colframe=#2,colupper=#3,
    fontupper=\bfseries,
    halign=center,valign=center,
    square,arc is angular,
    boxrule=0.8pt,boxsep=0.6mm,
}

\newtcbox{\qenumi}{enhanced,on line,size=minimal,boxsep=1pt,
    halign=center,valign=center,
    colback=qlabelbgcolor,colframe=qlabelolcolor,
    fontupper=\bfseries,colupper=qlabelfgcolor,
    boxrule=0.8pt,arc=2pt,outer arc=2pt,
    drop fuzzy shadow
    }

\newtcbox{\qenumii}{enhanced,on line,
    arc=0pt,outer arc=0pt,
    halign=center,valign=center,
    colback=qlabelbgcolor,colframe=qlabelolcolor,
    fontupper=\bfseries,colupper=qlabelfgcolor,
    boxsep=0pt,left=3pt,right=3pt,top=3pt,bottom=3pt,
    boxrule=0.0pt}





% from teaching
\NewTColorBox{skeleton}{ !O{} }{%
	enhanced jigsaw,colback=white,colframe=black,sharp corners,
    boxrule=0.6pt,
    fonttitle=\bfseries,colbacktitle=white,coltitle=black,
    boxed title style={left=0pt,colback=white,colframe=white},
    #1
}

\NewDocumentEnvironment{df}{ O{} }%
{\begin{skeleton}[title={Définition},colframe=dfcolor,#1]}{\end{skeleton}}

\NewDocumentEnvironment{cc}{ O{} }%
{\begin{skeleton}[title={Conclusion},colframe=cccolor,#1]}{\end{skeleton}}

\NewDocumentEnvironment{exe}{ O{} }%
{\begin{skeleton}[title={Exemple},colframe=execolor,#1]}{\end{skeleton}}

\NewDocumentEnvironment{rem}{ O{} }%
{\begin{skeleton}[title={Remarque},colframe=remcolor,#1]}{\end{skeleton}}

\newcounter{act}
\NewDocumentEnvironment{act}{ m }%
{%
    \begin{skeleton}[title={{\scshape Activité}~\refstepcounter{act}\theact~:~#1},
            breakable]
}%
{\end{skeleton}}

\newcounter{app}
\NewDocumentEnvironment{app}{ O{} }%
{%
    \begin{skeleton}[title={Application~\refstepcounter{app}\theapp},#1]
}%
{\end{skeleton}}

\NewTcbTheorem[number within=section]{Document}{Document}{docstyle}{th}


\NewTColorBox{Block}{ O{} }{enhanced jigsaw,
    colframe=black,opacityback=0.0,opacitybacktitle=0.0,
    boxrule=0.8pt,arc=0pt,outer arc=0.4pt,
    colbacktitle=white,coltitle=black,fonttitle=\bfseries\scshape,
    titlerule=0pt,top=2pt,bottom=3pt,left=2mm,right=2mm,
    breakable,
    title={#1},
}

%\NewDocumentCommand{\mathbox}{ m }%
%    {\tcboxmath[colframe=black,colback=boxbg,boxrule=0.6pt,sharp corners,
%                left=2mm,right=2mm,top=2mm,bottom=2mm]{#1}}


% from phys
% tcolorbox staff --------------------------------------------------------------
\colorlet{mathboxcolframe}{black}
\colorlet{mathboxcolback}{white}
\NewDocumentCommand{\mathbox}{ O{} O{mathboxcolframe} O{mathboxcolback} m }{%
    \tcboxmath[enhanced jigsaw,colframe=#2,colback=#3,opacityback=0.0,
        boxrule=0.8pt,arc=0pt,outer arc=0.4pt,
        left=0pt,right=0pt,top=0pt,bottom=0pt,
        boxsep=2mm,natural height,#1]{#4}
}


\NewTColorBox{donnees}{ !O{} }{skin=empty,
    title={Données~:},fonttitle=\bfseries\scshape,coltitle=black,
    boxrule=0pt,boxsep=0pt,left=0pt,right=0pt,top=2mm,
    before skip=2mm,after skip=2mm,
    before upper*={
        \renewcommand\labelitemi{\ding{121}}
        \renewcommand\labelitemii{$\bullet$}
    },
    #1
}

\NewTotalTCBox{\V}{ !O{} }
{fontupper=\bfseries,nobeforeafter,tcbox raise base,arc=2pt,%outer arc=0pt,
top=0pt,bottom=0pt,left=0mm,right=0mm,boxsep=0.5mm,
colback=Green!10,colframe=black,right skip=1pt,#1}{V}

\NewTotalTCBox{\F}{ !O{} }
{fontupper=\bfseries,nobeforeafter,tcbox raise base,arc=2pt,%outer arc=0pt,
top=0pt,bottom=0pt,left=0mm,right=0mm,boxsep=0.5mm,
colback=red!10,colframe=black,#1}{F}

\NewTColorBox{corrbox}{ !O{} }{
    skin=enhanced,,nobeforeafter,box align=top,
    sidebyside,righthand width=2cm,sidebyside align=center,
    % Other options: sidebyside gap=5mm, lefthand ratio=0.25, sidebyside adapt=right,
    segmentation style={solid,thick},
    #1
}

\NewDocumentEnvironment{eqdef}{ m }{%
    \begin{tcbitemize}[raster columns=2,raster equal height,
        raster force size=false,
        raster every box/.style={
            %size=small,
        },
        raster column 1/.style={halign=center,valign=center,
            add to width=-#1\textwidth,fontupper=\Large},
        raster column 2/.style={add to width=#1\textwidth}]
}{\end{tcbitemize}}
